\documentclass[10pt]{article}
\usepackage[a4paper, left=0.40in, right=0.40in, top=0.25in, bottom=0.15in]{geometry}
\usepackage[utf8]{inputenc}
\usepackage[english]{babel}
\hyphenation{Some-long-word}

\usepackage{resume}

\begin{document}
	
\begin{center}
	\textbf{\LARGE Vimarsh Sathia}\\[0.5ex]
	\textscale{1.05}{$1202$ Chelsea, Hiranandani Estate, Thane $400607$\\
		\Letter\hspace{0.1ex} \href{mailto:vimarsh.sathia@gmail.com}{vimarsh.sathia@gmail.com}
		\hfill
		\Mobilefone\hspace{0.1ex} +91 98194-33877}
\end{center}

\spacedhrule{-1.0ex}{-0.5ex}
\roottitle{EDUCATION}

	\headedsection
	{Indian Institute of Technology, Madras}
	{Chennai, India}
	{Bachelors of Technology in Computer Science}
	{\period{July}{2017}{June}{2021(expected)}}
	{GPA : $9.1$/$10$ (after $6$ semesters)}

\spacedhrule{0.8ex}{0.0ex}
\roottitle{ACHIEVEMENTS}

\begin{circlist}
	\item Awarded the first place at JARVIS - the Data Analytics hackathon event at \href{https://en.wikipedia.org/wiki/Shaastra}{Shaastra 2020}, an All India tech fest organized by IIT Madras
	\item Awarded a branch change from Engineering Physics to Computer Science on the basis of GPA at the end of Semester $1$
	\item Secured an All India Rank of $1806$ among half a million candidates in \href{https://en.wikipedia.org/wiki/Joint_Entrance_Examination_%E2%80%93_Advanced}{JEE} $2017$.
\end{circlist}

\spacedhrule{0.8ex}{0.0ex}
\roottitle{WORK EXPERIENCE}
	
\headedsection
{\href{https://www.microsoft.com/en-in/}{Microsoft}}
{Hyderabad, India}
{Software Engineering Internship}
{\period{May}{2020}{July}{2020}}
{
	\vspace{-2.4ex}
	\begin{circlist}
		\item Designed an API for bulk reading of contacts in a database, with support for asynchronous scheduling and execution based on the availability of server resources.
		\item Tested and delivered a succesful end to end implementation of the same in C\#
		\item Performed further optimizations on the API to reduce the intermediate metadata size used from $400$ MB to $20$ MB.
	\end{circlist}
}

\headedsection
{\href{https://actifydatalabs.com/}{Actify Data Labs}}
{Bangalore, India}
{Analytics Intern}
{\period{June}{2019}{July}{2019}}
{
	\vspace{-2.4ex}
	\begin{circlist}
		\item Worked on generating synthetic data from vector training datasets
		\item Implemented the above in Java by building empirical copula models and drawing output points from them.
	\end{circlist}
}

\headedsection
{\href{https://www.dxc.technology/}{DXC Technology}}
{Bangalore, India}
{Intern}
{\period{June}{2018}{July}{2018}}
{
	\vspace{-2.4ex}
	\begin{circlist}
		\item Implemented standard quantum computing algorithms like Grover's Search and Shor's Prime Factoring using the Q\# language. Featured in company tech blog posts \href{https://blogs.dxc.technology/2018/07/30/what-i-learned-quantum-computing-experiment/}{here} and \href{https://blogs.dxc.technology/2018/08/08/my-experience-as-a-dxc-labs-intern/}{here}. 
	\end{circlist}
}

\spacedhrule{0.8ex}{0.0ex}
\roottitle{PROJECTS}
	
	\headedsectionthree
	{Conway's Game of Life}
	{\period{Jan}{2020}{May}{2020}}
	{CS6023: GPU Programming}
	{\vspace{-2.4ex}
		\begin{circlist}
			\item Course project on a parallel implementation and visualization of cellular automation using GPUs.
			\item Implemented parallel computation using the Nvidia CUDA library, and rendered the results using texture mapping in OpenGL. 
		\end{circlist}
	}

	\headedsectionthree
	{Hidden Markov Models for Spoken Digit and Handwritten Character Classification}
	{\period{Jul}{2019}{Nov}{2019}}
	{CS5691: Pattern Recognition and Machine Learning}
	{\vspace{-2.4ex}
		\begin{circlist}
			\item Modelled spoken digits and handwritten Telugu characters using left to right HMM sequences trained using the Baum-Welch algorithm for classification
			\item Achieved multi-digit and multi-character classification by concatenating individually trained HMM sequences.
		\end{circlist}
	}
	
	\headedsectionthree
	{MiniJava Compiler}
	{\period{Jul}{2019}{Nov}{2019}}
	{CS3300: Compilers}
	{\vspace{-2.4ex}
		\begin{circlist}
			\item Course assignment on writing a $6$ pass compiler for minijava (a subset of the java language) and converting the code to MIPS assembly.
			\item Implemented using \href{https://javacc.github.io/javacc/}{JavaCC} as the parser generator, with \href{http://compilers.cs.ucla.edu/jtb/}{JTB} as a syntax tree builder for all intermediate passes.
		\end{circlist}
	}

	\headedsectionthree
	{Prolog interpreter in OCaml}
	{\period{Jul}{2019}{Nov}{2019}}
	{CS3100: Paradigms of Programming}
	{\vspace{-2.4ex}
		\begin{circlist}
			\item Final course assignment to implement a Prolog interpreter in OCaml with support for backtracking(recover from bad choices) and choice points(produce multiple results).
		\end{circlist}
	}

	% \headedsectionthree
	% {HTTP Proxy Server}
	% {\period{Jan}{2020}{May}{2020}}
	% {CS3205: Computer Networks}
	% {\vspace{-2.4ex}
	% 	\begin{circlist}
	% 		\item Implemented a http proxy server in C, with support for concurrent and split http GET requests using the HTTP/1.0 specification
	% 	\end{circlist}
	% }

\spacedhrule{0.8ex}{0.0ex}
\roottitle{RELEVANT COURSEWORK \& SKILLS}

	\begin{indentsection}
		\skill{Machine Learning}{Pattern Recognition and Machine Learning, Numerical Optimization, Applied Statistics, Deep Learning}
		\skill{Systems}{Operating Systems, Computer Networks, GPU Programming, Blockchain and Distributed Ledger Tech}
		\skill{PL}{Compilers, Paradigms of Programming(Functional Programming)}
		\skill{Algorithms}{Design and Analysis of Algorithms, Topics in Algorithm Design, Pseudorandomness}
	\end{indentsection}

\spacedhrule{0.8ex}{0.0ex}
\roottitle{EXTRA-CURRICULARS}
	\begin{indentsection}
		\skill{Squash}{Member of the institute squash team at IIT Madras }
		\skill{Comedy}{Member of the Comedy Club at IIT Madras, participating and organizing open mics and improv meetups}
	\end{indentsection}

\end{document}
